\nonstopmode{}
\documentclass[letterpaper]{book}
\usepackage[times,inconsolata,hyper]{Rd}
\usepackage{makeidx}
\makeatletter\@ifl@t@r\fmtversion{2018/04/01}{}{\usepackage[utf8]{inputenc}}\makeatother
% \usepackage{graphicx} % @USE GRAPHICX@
\makeindex{}
\begin{document}
\chapter*{}
\begin{center}
{\textbf{\huge Package `SII'}}
\par\bigskip{\large \today}
\end{center}
\ifthenelse{\boolean{Rd@use@hyper}}{\hypersetup{pdftitle = {SII: Calculate ANSI S3.5-1997 Speech Intelligibility Index}}}{}
\ifthenelse{\boolean{Rd@use@hyper}}{\hypersetup{pdfauthor = {Gregory R. Warnes; Mark Shaver}}}{}
\begin{description}
\raggedright{}
\item[Type]\AsIs{Package}
\item[Title]\AsIs{Calculate ANSI S3.5-1997 Speech Intelligibility Index}
\item[Version]\AsIs{1.1.0}
\item[Date]\AsIs{2026-07-19}
\item[Description]\AsIs{
Calculates ANSI S3.5-1997 Speech Intelligibility Index
('SII'), a standard method for computing the intelligibility of
speech from acoustical measurements of speech, noise, and hearing
thresholds. This package includes data frames corresponding to
Tables 1 - 4 in the ANSI standard as well as a function utilizing
these tables and user-provided hearing threshold and noise level
measurements to compute the SII score.  The methods implemented
here extend the standard computations to allow calculation of SII
when the measured frequencies do not match those required by the
standard by applying interpolation to obtain values for the
required frequencies
-- 
Development of this package was funded by the Center for 'Bioscience'
Education and Technology ('CBET') of the Rochester Institute of
Technology ('RIT').}
\item[Suggests]\AsIs{splines, readxl, xtable, gdata, shiny, bslib}
\item[License]\AsIs{GPL-2}
\item[LazyLoad]\AsIs{yes}
\item[NeedsCompilation]\AsIs{no}
\item[Repository]\AsIs{CRAN}
\item[Date/Publication]\AsIs{2018-11-18 09:56:48 UTC}
\item[Author]\AsIs{Gregory R. Warnes [aut],
Mark Shaver [cre, aut]}
\item[Maintainer]\AsIs{Mark Shaver }\email{mark.shaver@posteo.net}\AsIs{}
\end{description}
\Rdcontents{Contents}
\HeaderA{SII-package}{Calculate ANSI S3.5-1997 Speech Intelligibility Index}{SII.Rdash.package}
\aliasA{SII}{SII-package}{SII}
\keyword{package}{SII-package}
%
\begin{Description}
This package calculates ANSI S3.5-1997 Speech Intelligibility Index
(SII), a standard method for computing the intelligibility of
speech from acoustical measurements of speech, noise, and hearing
thresholds. This package includes data frames corresponding to
Tables 1 - 4 in the ANSI standard as well as a function utilizing
these tables and user-provided hearing threshold and noise level
measurements to compute the SII score.  The methods implemented
here extend the standard computations to allow calculation of SII
when the measured frequencies do not match those required by the
standard by applying interpolation to obtain values for the
required frequencies.
\end{Description}
%
\begin{Author}
Gregory R. Warnes \email{greg@warnes.net}
\end{Author}
%
\begin{References}
ANSI S3.5-1997, "American National Standard Methods for 
Calculation of the Speech Intelligibility Index" American National Stan- 
dards Institute, New York.

Other software programs for calculating SII are available from
\url{https://sii.to/html/programs.html}.
\end{References}
%
\begin{Examples}
\begin{ExampleCode}
## Example C.1 from ANSI S3.5-1997 Annex C
sii.C1 <- sii(
              speech   = c(50.0, 40.0, 40.0, 30.0, 20.0,  0.0),
              noise    = c(70.0, 65.0, 45.0, 25.0,  1.0,-15.0),
              threshold= c( 0.0,  0.0,  0.0,  0.0,  0.0,  0.0),
              method="octave"
	      )
sii.C1                        # rounded to 2 digits by default
print(sii.C1$sii, digits=20)  # full precision
summary(sii.C1)               # full details
plot(sii.C1)                  # plot
## The value given in the Standard is $0.504$.
\end{ExampleCode}
\end{Examples}
\HeaderA{calculate\_loudness}{Calculate Loudness (Sones)}{calculate.Rul.loudness}
%
\begin{Description}
Calculates the estimated loudness in sones for a given Speech Intelligibility Index (SII) object.
\end{Description}
%
\begin{Usage}
\begin{verbatim}
calculate_loudness(x)
\end{verbatim}
\end{Usage}
%
\begin{Arguments}
\begin{ldescription}
\item[\code{x}] An object of class \code{SII}, representing the output of the \code{sii()} function.
\end{ldescription}
\end{Arguments}
%
\begin{Value}
Returns a numeric value representing the total loudness in sones.
\end{Value}
%
\begin{Author}
Mark Shaver
\end{Author}
%
\begin{Examples}
\begin{ExampleCode}
sii.res <- sii(speech=c(50, 40, 40, 30, 20, 0), threshold=rep(0,6), method="octave")
calculate_loudness(sii.res)
\end{ExampleCode}
\end{Examples}
\HeaderA{critical}{Constants Tables for ANSI S3.5-1997 Speech Intelligibility Index (SII)}{critical}
\aliasA{equal}{critical}{equal}
\aliasA{octave}{critical}{octave}
\aliasA{onethird}{critical}{onethird}
\aliasA{overall.spl}{critical}{overall.spl}
\keyword{datasets}{critical}
%
\begin{Description}
Tables of constants for ANSI S3.5-1997 Speech Intelligibility Index
(SII) 
\end{Description}
%
\begin{Usage}
\begin{verbatim}
data(critical)
data(equal)
data(onethird)
data(octave)
data(overall.spl)
\end{verbatim}
\end{Usage}
%
\begin{Format}
Each data frames has 6-21 observations and a subset of the following
variables: 
\begin{description}

\item[\code{fi}] Center frequency of SII band, Hz
\item[\code{li}] Lower limit of frequency band, Hz
\item[\code{hi}] Upper limit of frequency band, Hz
\item[\code{Deltai}] Band width adjustment, dB
\item[\code{Ii}] Band importance function
\item[\code{normal}, \code{raised}, \code{loud} and
\code{shout}] Standard spectrum levels for vocal effort levels
"normal", "raised", "loud", and "shout", respectively, dB
\item[\code{Xi}] Spectrum level of internal noise, dB
\item[\code{Fi}] Band importance function (weight)

\end{description}

\end{Format}
%
\begin{Details}
These data objects provide constant tables 1 -- 4 from the ANSI
S3.5-1997.
\begin{description}

\item[\code{critical}] Table 1: Critical band SII procedure
constants
\item[\code{equal}] Table 2:Equally contributing (17 band) critical
band SII
\item[\code{onethird}] Table 3: One-third octave band SII procedure
constants
\item[\code{octave}] Table 4: Octave band SII procedure constants
\item[\code{overall.spl}] Overall sound pressure level (SPL) for the
for  vocal effort levels "normal", "raised", "loud", and "shout",
in dB

\end{description}

\end{Details}
%
\begin{Source}
ANSI S3.5-1997, "American National Standard Methods for 
Calculation of the Speech Intelligibility Index" American National
Standards Institute, New York. 
\end{Source}
%
\begin{References}
ANSI S3.5-1997, "American National Standard Methods for 
Calculation of the Speech Intelligibility Index" American National
Standards Institute, New York. 
\end{References}
%
\begin{Examples}
\begin{ExampleCode}
data(critical)
critical # show entire table

data(equal)
names(equal)
equal$fi # extract just the frequency band centers

data(onethird)
barplot(onethird$Ii) # plot band importance function (weights)

data(octave)
round(octave, digits=2) # just 2 digits

data(overall.spl)
overall.spl
\end{ExampleCode}
\end{Examples}
\HeaderA{launch\_app}{Launch the SII Interactive Dashboard}{launch.Rul.app}
\keyword{misc}{launch\_app}
%
\begin{Description}
Launches a local Shiny web application that provides an interactive graphical interface to the SII calculation engine and plotting utilities.
\end{Description}
%
\begin{Usage}
\begin{verbatim}
launch_app()
\end{verbatim}
\end{Usage}
%
\begin{Details}
The interactive dashboard allows you to dynamically adjust hearing thresholds across standard clinical frequencies (250 Hz - 8000 Hz) using sliders. It provides real-time visualization of the Clinical SPLogram and the Insertion Gain / Compression plot based on the selected prescriptive rationale (e.g., Unaided, NAL-R, Open-NL) and desensitization settings.

Note: This function requires the \code{shiny} and \code{bslib} packages to be installed.
\end{Details}
%
\begin{Author}
Gregory R. Warnes \email{greg@warnes.net}, Mark Shaver \email{mark.shaver@posteo.net}
\end{Author}
%
\begin{SeeAlso}
\code{\LinkA{sii}{sii}}, \code{\LinkA{plot.SII}{plot.SII}}
\end{SeeAlso}
%
\begin{Examples}
\begin{ExampleCode}
## Not run: 
  # Launch the interactive SII dashboard in your web browser
  launch_app()

## End(Not run)
\end{ExampleCode}
\end{Examples}
\HeaderA{predict\_aided\_sii}{Predict Aided Speech Intelligibility Index (SII)}{predict.Rul.aided.Rul.sii}
%
\begin{Description}
Uses multiple linear regression equations from Johnson (2013) to accurately predict 
the aided Speech Intelligibility Index (SII) for a 65 dB SPL average speech input. 
The predictions approximate the full 2,160-step ANSI S3.5 calculation based on pure 
tone averages (PTA).
\end{Description}
%
\begin{Usage}
\begin{verbatim}
predict_aided_sii(freq, threshold, age = c("adult", "pediatric"), prescription = c("NAL-NL2", "DSL"), desensitized = FALSE)
\end{verbatim}
\end{Usage}
%
\begin{Arguments}
\begin{ldescription}
\item[\code{freq}] A numeric vector of frequencies (in Hz) corresponding to the \code{threshold} values.
\item[\code{threshold}] A numeric vector of hearing thresholds (in dB HL).
\item[\code{age}] A character string specifying the patient age group. Must be either \code{"adult"} (default) or \code{"pediatric"} (modeled on 3-year-olds).
\item[\code{prescription}] A character string specifying the prescriptive rationale to predict. Must be either \code{"NAL-NL2"} (default) or \code{"DSL"} (DSL m[i/o] v5.0).
\item[\code{desensitized}] Logical flag. If \code{TRUE}, predicts the effective SII incorporating hearing loss desensitization. If \code{FALSE} (default), predicts the traditional ANSI SII.
\end{ldescription}
\end{Arguments}
%
\begin{Value}
A numeric value representing the predicted SII (constrained between 0.0 and 1.0).
\end{Value}
%
\begin{Author}
Maintainer
\end{Author}
%
\begin{References}
Johnson, E. E. (2013). Modern prescription theory and application: realistic expectations for speech recognition with hearing aids. Trends in Amplification, 17(3/4), 143-170.
\end{References}
\HeaderA{sic.critical}{Alternative ANSI S3.5-1997 SII Transfer Function Weights}{sic.critical}
\aliasA{sic.octave}{sic.critical}{sic.octave}
\aliasA{sic.onethird}{sic.critical}{sic.onethird}
\keyword{datasets}{sic.critical}
%
\begin{Description}
Alternative ANSI S3.5-1997 Speech Intelligibility Index (SII) transfer
function weights for for various types of speech material.
\end{Description}
%
\begin{Usage}
\begin{verbatim}
data(sic.critical)
data(sic.onethird)
data(sic.octave)
\end{verbatim}
\end{Usage}
%
\begin{Format}
Each data frame  contains the following 8 variables, each
corresponding the the transfer function weights for a specific type of
speech material:
\begin{description}

\item[\code{fi}] Center frequency, Hz
\item[\code{SII}] Standard SII transfer function (weights)
\item[\code{NNS}] NNS (various nonsense syllable tests where
most of the English phonems occur equally often)
\item[\code{CID22}] CID-W22 (PB-words)
\item[\code{NU6}] NU6 monosyllables
\item[\code{DRT}] DRT (Diagnostic Rhyme Test)
\item[\code{ShortPassage}] short passages of easy reading material
\item[\code{SPIN}] SPIN monosyllables
\item[\code{CST}] Connected Speech Test

\end{description}

\end{Format}
%
\begin{Details}
\begin{description}

\item[\code{sic.critical}] provides alternative weights for the
critical band SII procedure.
\item[\code{sic.threeoctave}] provides alternative weights for the
one-third octave frequency band SII procedure.
\item[\code{octave}] provides alternative weights for the
octave frequency band SII procedure.

\end{description}

\end{Details}
%
\begin{Section}{note}
There is no table of alternative weights for the
equally-weighted SII band procedure as the weights for this method are
(by definition) constant across all bands.

\end{Section}
%
\begin{Source}
All values except the \code{CST} columns are from:

ANSI S3.5-1997, "American National Standard Methods for 
Calculation of the Speech Intelligibility Index" American National
Standards Institute, New York. 

Values in the \code{CST} columns are from the original Connected Speech Test (CST) dataset.
\end{Source}
%
\begin{References}
ANSI S3.5-1997, "American National Standard Methods for 
Calculation of the Speech Intelligibility Index" American National
Standards Institute, New York. 
\end{References}
%
\begin{Examples}
\begin{ExampleCode}
## Load the alternative weights for the critical band method
data(sic.critical)

## display the weights
round(sic.critical,3)

## draw a comparison plot
ngroup <- ncol(sic.critical)
matplot(x=sic.critical[,1], y=sic.critical[,-1],
        type="o",
        xlab="Frequency, Hz",
        ylab="Weight",
        log="x",
        lty=1:ngroup,
        col=rainbow(ngroup)
)
legend(
       "topright",
       legend=names(sic.critical)[-1],
       pch=as.character(1:ngroup),
       lty=1:ngroup,
       col=rainbow(ngroup)
       )

data(threeoctave)
data(octave)
\end{ExampleCode}
\end{Examples}
\HeaderA{sii}{Compute ANSI S3.5-1997 Speech Intelligibility Index (SII)}{sii}
\aliasA{plot.SII}{sii}{plot.SII}
\aliasA{print.SII}{sii}{print.SII}
\aliasA{summary.SII}{sii}{summary.SII}
\keyword{math}{sii}
%
\begin{Description}
Compute the Speech Intelligibility Index (SII) described by ANSI
specification S3.5-1997, including extensions for conductive
hearing loss. Optionally apply interpolation obtain values for the
required frequencies.
\end{Description}
%
\begin{Usage}
\begin{verbatim}
sii(speech = c("normal", "raised", "loud", "shout"),
    noise, threshold, loss, freq, 
    method = c("critical", "equal-contributing",
               "one-third octave", "octave"),
    importance = c("SII", "NNS", "CID22", "NU6", "DRT",
                   "ShortPassage", "SPIN", "CST"),
    interpolate=FALSE,
    prescription=NULL,
    desensitization=FALSE,
    ldl=NULL,
    gender="male",
    experience="experienced",
    config="bilateral",
    age="adult",
    age_years=NULL,
    age_months=NULL,
    coupling="custom_occluded",
    module="standard",
    transducer="inserts",
    custom_gain=NULL)
## S3 method for class 'SII'
print(x, digits=3, ...)
## S3 method for class 'SII'
plot(x, clinical=FALSE, legend=TRUE, legend_only=FALSE, ...)
## S3 method for class 'SII'
summary(object, digits=2, ...)
\end{verbatim}
\end{Usage}
%
\begin{Arguments}
\begin{ldescription}
\item[\code{speech}] Either a numeric vector providing \eqn{E'_i}{}, the
equivalent speech spectrum level (in dB) at each frequency, or a
character string indicating the stated vocal effort corresponding to
one of the standard standard speech spectrum levels ("normal",
"raised", "loud", "shout"). Defaults to \code{speech="normal"}
correspoding to the normal level of stated vocal effort. 
\item[\code{noise}] A numeric vector providing \eqn{N'_i}{}, the
equivalent noise spectrum level (in dB) at each frequency.  If
missing, defaults to -50 dB for each frequency.
\item[\code{threshold}] A numeric vector providing \eqn{T'_i}{}, the
equivalent hearing threshold level (in dB) at each frequency.  If
missing, defaults to 0 dB for each frequency.
\item[\code{loss}] A numeric vector providing \eqn{J'_i}{}, the
conductive hearing loss level (in dB) at each frequency.  If
missing, defaults to 0 dB for each frequency.
\item[\code{freq}] Vector of frequencies for which \code{speech},
\code{noise}, \code{threshold}, and/or \code{loss} are specified.
If \code{interpolate=TRUE}, \code{freq} must be specified.
Otherwise, it must either match the required value for SII
calculation method given by argument \code{method}, or be missing,
in which case it will default to the values required for the
specified method.
\item[\code{method}] A character string specifying the SII calculation
method ("critical", "one-third octave", "equal-contributing",
"octave")
\item[\code{importance}] Either a numeric vector providing \eqn{F_i}{},
the transfer function (importance weights) at each frequency, or a
character string indicating which transfer function to employ
("SII", "NNS", "CID22", "NU6", "DRT", "ShortPassage", "SPIN",
"CST"). Defaults to the standard SII transfer function,
\code{importance="SII"}.
\item[\code{interpolate}] Logical flag indicating whether to interpolate 
from the provide measurement values and frequencies to those 
required by the specified method via linear interpolation on the log
scale.
\item[\code{prescription}] A character string (e.g. \code{"NAL-R"}, \code{"Open-NL"}). If provided, the function will automatically calculate hearing aid insertion gain based on the specified prescriptive fitting rationale and add it to the speech and noise spectrum, calculating an Aided SII.
\item[\code{desensitization}] Logical flag. If \code{TRUE}, applies the hearing loss desensitization factor (Ching et al., 2011; Johnson, 2013) to the calculation to account for suprathreshold distortion in impaired ears. Defaults to \code{FALSE} (Traditional ANSI S3.5 calculation).
\item[\code{ldl}] Numeric vector specifying Loudness Discomfort Levels (LDL) in dB HL for each frequency. If \code{NULL}, limits are estimated based on hearing thresholds.
\item[\code{gender}] Character string specifying patient gender ("male" or "female"). Used for prescriptive algorithms. Defaults to "male".
\item[\code{experience}] Character string specifying user hearing aid experience ("experienced", "new", "power"). Defaults to "experienced".
\item[\code{config}] Character string for fitting configuration ("bilateral" or "unilateral"). Defaults to "bilateral".
\item[\code{age}] Character string for patient age group (e.g., "adult", "child\_0\_5", "child\_6\_11", "child\_12\_23", "child\_24\_35", "child\_36\_59"). Used for RECD corrections and pediatric prescriptive modifications.
\item[\code{age\_years}] Optional. Integer specifying exact age in years, overriding the categorical `age` bin if provided.
\item[\code{age\_months}] Optional. Integer specifying exact age in months. Only applies to pediatric RECD scaling.
\item[\code{coupling}] Character string specifying the acoustic coupling / vent ("custom\_occluded", "double\_dome", "tulip\_dome", "open\_dome", "vent\_1mm\_solid", "vent\_2mm\_solid", "vent\_3mm\_solid", "vent\_1mm\_hollow", "vent\_2mm\_hollow", "vent\_3mm\_hollow"). Modifies low-frequency leakage and insertion gain targets.
\item[\code{module}] Character string for hearing aid module ("standard" or other specific module constraints).
\item[\code{transducer}] Character string specifying the audiometric transducer used for testing ("inserts" or "supra\_aural"). Determines RETSPL and RECD corrections for SPLogram plots.
\item[\code{custom\_gain}] Numeric vector specifying a custom insertion gain array (in dB) for each frequency. If provided along with \code{prescription="Custom"}, this gain is explicitly applied.
\item[\code{object}, \code{x}] SII object
\item[\code{digits}] Number of digits to display
\item[\code{clinical}] Logical flag. If \code{TRUE}, plots a patient-friendly clinical SPLogram highlighting the audible speech area. If \code{FALSE} (the default), the behavior depends on whether the object is aided or unaided. For unaided objects, it plots a diagnostic interpolation graph. For aided objects (when a prescription was used), it automatically generates a 3-line Insertion Gain plot for 55, 65, and 75 dB SPL input levels (dynamically loading and scaling the true Normal and Loud LTASS spectra).
\item[\code{legend}] Logical flag. If \code{TRUE} (the default), draws a legend on the plot.
\item[\code{legend\_only}] Logical flag. If \code{TRUE}, suppresses the main plot and only draws the legend in the center of the plotting window. Useful for multi-plot layouts.
\item[\code{...}] Optional arguments to \code{print}, \code{summary}, and
\code{plot} methods
\end{ldescription}
\end{Arguments}
%
\begin{Details}
American National Standard ANSI S3.5-1997 ("Methods for Calculation of
the Speech Intelligibility Index") defines a method for computing a
physical measure that is highly correlated with the intelligibility
of speech as evaluated by speech perception tests given a group of
talkers and listeners. This measure is called the Speech
Intelligibility Index, or SII. The SII is calculated from acoustical
measurements of speech and noise.

The \code{sii} function implements ANSI S3.5-1997 as described in
the standard, without any attempt to optimize the performance. The
implementation does, however, include the extension for handling
conductive hearing loss from Annex A (utilizing the optional
\code{loss} argument), and for utilizing alternative band weights
(i.e. transfer function) appropriate for differing message contents
(e.g. types of speech) as described in Annex B or user-specified
band weights (utilizing the optional argument \code{importance}).

Further, this implementation provides a mechanism for
interpolating/extrapolating available measurements to those required
for the specified calculation procedure.  When
\code{interpolate=TRUE}, required values for \code{speech},
\code{noise}, \code{threshold}, and \code{loss} will be computed 
using linear interpolation (of the log-scaled data).  In this case,
missing values may be provided and will be appropriately
interpolated.

%
\begin{SubSection}{Prescriptive Fitting Rationales}
If the \code{prescription} argument is provided, the function will dynamically calculate a frequency-specific hearing aid insertion gain and apply it to the speech and noise spectrum, calculating an Aided SII. The following rationales are supported:

\begin{itemize}

\item{} \code{"NAL-R"}: A classic linear fitting rationale designed by the National Acoustic Laboratories to maximize speech intelligibility for mild-to-moderate losses. Gain is a linear function of the pure-tone average (PTA) and frequency-specific thresholds.
\item{} \code{"Open-NL"}: A completely transparent, open-source non-linear fitting algorithm developed specifically for this package. It was designed to serve as an open-source alternative to proprietary modern non-linear clinical targets (such as NAL-NL2 and DSL v5.0). The algorithm calculates insertion gain using the following explicit mathematical steps:
\begin{enumerate}

\item{} \strong{Conductive Component Separation}: If an Air-Bone Gap is present, the purely sensorineural component is isolated for WDRC compression. Linear gain representing 75\% of the conductive loss is added at the end (consistent with NAL-NL2/Johnson).
\item{} \strong{Dynamic Base Gain}: The base gain relies on a dynamic multiplier (ranging from 0.43 to 0.48) determined by user experience, combined with empirical NAL-R shaping constants. Severe losses receive a booster, while steep slopes receive a low-frequency penalty to avoid upward spread of masking.
\item{} \strong{Dead Region Roll-offs}: High-frequency or low-frequency dead regions trigger a steep 30 dB/octave penalty beyond the viable boundaries to prevent acoustic distortion and feedback.
\item{} \strong{Wide Dynamic Range Compression (WDRC)}: A dynamic compression ratio (ranging from 1.0 to 3.0) is applied symmetrically around a 65 dB SPL pivot point. The ratio scales dynamically based on the sensorineural threshold and measured Loudness Discomfort Levels (LDLs).
\item{} \strong{Bandwidth Roll-off}: Empirical bandwidth roll-offs are applied to limit unnecessary low/high frequency amplification, with separate curves for adults vs. infants.
\item{} \strong{Acoustic Coupling}: Adjustments are applied to the low frequencies depending on the specified venting, dome type, and acoustic seal.
\item{} \strong{MPO Limits}: Output is capped by predictive NAL-SSPL90 limits, adjusting for conductive attenuation, with an absolute safety hard cap at 120 dB SPL at the cochlea and a hardware output limit of 135 dB SPL.

\end{enumerate}


\end{itemize}


\end{SubSection}

\end{Details}
%
\begin{Value}
The return value is an object of class SII, containing the following
components:
\begin{ldescription}
\item[\code{call}] Function call used to generate the SII object
\item[\code{orig}] List containing original (pre-extrapolation) values for
\code{freq}, \code{speech}, \code{noise}, \code{threshold}, and
\code{loss}.
\item[\code{speech}, \code{noise}, \code{threshold}, \code{loss}, \code{and freq}] Values used in
calculations (extrapolated if necessary)
\item[\code{unaided\_speech}] Original speech array before applying any prescriptive gain
\item[\code{vocal\_effort}] String representing the stated vocal effort
\item[\code{gain}] Insertion gain array added to speech and noise when a prescription is used
\item[\code{prescription}] The fitting rationale used (if any)
\item[\code{unaided\_sii}] Calculated SII value before applying the prescriptive gain (if a prescription was used)
\item[\code{table}] SII calculation worksheet, containing columns
corresponding to both Table C.1 and C.2 in Annex C of the
standard. Table columns are
\begin{description}

\item[Fi] Center frequency of SII band, Hz
\item[E'i] Spectrum level of equivalent speech, dB
\item[N'i] Spectrum level of equivalent noise, dB
\item[T'i] Equivalent hearing threshold level, dB
\item[Vi] Spectrum level for self-speech masking, dB
\item[Bi] Larger of the specrum levels for equivalent noise and
self-speech masking, dB
\item[Ci] Slope per octave (doubling of frequency) of the upward
spread of masking, dB/octave
\item[Zi] Spectrum level for equivalent masking, dB
\item[Xi] Spectrum level of internal noise, dB
\item[X'i] Spectrum level of equivalent internal noise, dB
\item[Di] Spectrum level for equivalent disturbance, dB
\item[Ui] Spectrum level of standard speech for normal vocal
effort, dB
\item[Ji] Equivalent hearing threshold due to conductive hearing
loss, dB
\item[Li] Speech level distortion factor, dB
\item[Ki] Temporary variable used in the calculation of the band
auditability function
\item[Ai] Band auditability function
\item[Ii] Band importance function
\item[IiAi] Product of the band importance function (Ii), and band
auditability function(Ai)

\end{description}


\item[\code{sii}] Calculated SII value
\end{ldescription}
\end{Value}
%
\begin{Section}{Open-NL Prescription Rationale}
\strong{DISCLAIMER:} \code{Open-NL} is an untested, experimental fitting rationale. It is designed to mimic aspects of other generic non-linear fitting rationales (such as NAL-NL2 and DSL v5.0) strictly to promote open research, algorithmic transparency, and rapid iteration within the audiology community. It is not intended for clinical use.

When \code{prescription = "Open-NL"} is selected, the function calculates and applies this WDRC-optimized rationale to maximize the aided Speech Intelligibility Index.

The mathematical algorithm consists of:
\begin{enumerate}

\item{} \strong{Minimal Hearing Loss (MHL) Bypass:} If \code{module == "mhl"} and \eqn{PTA_{.5,1,2,4k} \le 25}{} dB HL, WDRC is bypassed. Applies flat linear gain interpolated from \eqn{(f, G) = \{(250,0), (500,0), (1k,3), (2k,5), (4k,5), (8k,5)\}}{} dB, with a 1.5 compression ratio for loud inputs.
\item{} \strong{Base Anchor (65 dB SPL Input):} \eqn{g_{65} = \max(0, m \times HTL + C)}{}, where \eqn{m \in \{0.40, 0.45, 0.50\}}{} depending on \code{experience}, and \eqn{C}{} interpolates arrays such as \eqn{\{-8, -1, +3, +1, 0, 0, 0, 0\}}{} dB evaluated at standard audiometric frequencies.
\item{} \strong{Audiometric Profile Corrections:}
\begin{itemize}

\item{} \emph{Steep Slope Knee:} For slopes \eqn{>30}{} dB/octave, penalizes the "knee" (500-1500 Hz) by up to 6 dB, and boosts frequencies \eqn{\ge 2000}{} Hz by up to 6 dB.
\item{} \emph{Severe-Loss Booster (SLB):} Adds \eqn{\min(15, \max(0, HTL - 60) \times 0.5)}{} dB. Tapered heavily in mid-frequencies and disabled in dead regions.

\end{itemize}

\item{} \strong{High-Frequency Desensitization:} To prevent distortion, excess gain above a limit is compressed at a 2:1 ratio.
\begin{itemize}

\item{} Adult \eqn{Limit = 30 + 0.4 \times \max(0, HTL - 60)}{}
\item{} Child \eqn{Limit = 40 + 0.5 \times \max(0, HTL - 60)}{}

\end{itemize}

\item{} \strong{Roll-offs and Dead Regions:}
\begin{itemize}

\item{} \emph{Bandwidth Roll-off:} Multiplies \eqn{g_{65}}{}. Adults: interpolates \eqn{\{(250Hz, 0.7), (500, 1.0), ..., (6k, 0.8), (8k, 0.5)\}}{}. Children: 1.0 across all frequencies.
\item{} \emph{Dead Regions:} Identifies HF Dead Regions (\eqn{HTL \ge 90}{} at \eqn{f \ge 1k}{}) and LF Dead Regions (\eqn{HTL \ge 80}{} at \eqn{f \le 1k}{}). Applies a 30 dB/octave penalty beyond viable boundaries (\eqn{1.7 \times f_{e}}{} for HF, \eqn{0.57 \times f_{e}}{} for LF).

\end{itemize}

\item{} \strong{Bi-directional WDRC:}
\begin{itemize}

\item{} \emph{Compression Ratio (CR):} \eqn{CR_{base} = 1 + \max(0, HTL - 20) / 40}{}. For \eqn{HTL > 65}{}, CR reduces toward 1.0 in low frequencies. Overall bounded strictly to \eqn{\le 1.5}{} for low frequencies (\eqn{\le 500}{} Hz) and up to \eqn{2.4}{} for high frequencies (\eqn{\ge 3000}{} Hz).
\item{} \emph{Compression Threshold (CT):} Interpolated from \eqn{(HTL, CT)}{}. Adult: \eqn{\{(20HL, 30SPL), ..., (100, 45)\}}{}. Child: \eqn{\{(20, 25), ..., (100, 40)\}}{}.
\item{} \emph{Aggressive MPO Defense:} For steeply sloping losses (\eqn{>15}{} dB difference) with low LDLs (\eqn{<100}{} dB SPL), the formula proactively defends against MPO collision. The \eqn{CT}{} is aggressively lowered by up to 10 dB, and the \eqn{CR}{} is forced up by an additional 0.05 per dB of LDL penalty.
\item{} \emph{I/O Computation:} Calculates gain at CT (\eqn{G_{CT}}{}) scaling back from 65 dB SPL pivot. Applies linear gain below CT, and WDRC above CT.

\end{itemize}

\item{} \strong{Empirical Adjustments and Smoothing:}
\begin{itemize}

\item{} \emph{Demographic Boosts:} Child Boost: Interpolates \eqn{(HTL, Boost) = \{(20, 5), (50, 5), (80, 2), (100, 1)\}}{} dB. Gender: Female \eqn{= -1.5}{} dB. Config: Unilateral \eqn{= +3.0}{} dB. Experience: New users with \eqn{PTA > 40}{} get up to \eqn{-6.0}{} dB penalty.
\item{} \emph{LDL Dynamic Range Mapping:} If \code{ldl} is provided, dynamic range is evaluated. For every 1 dB the measured LDL is lower than predicted, \eqn{g_{65}}{} is reduced by 0.2 dB and \eqn{CR_{base}}{} is increased by 0.02.
\item{} \emph{Acoustic Coupling:} Subtracts vent leakage (e.g., Open Dome \eqn{= \{-35, -28, -15, -2, 0, 0\}}{} dB).
\item{} Applies a 3-point moving average to the final gain array.

\end{itemize}

\item{} \strong{SSPL90 MPO Limiting:} 
\begin{itemize}

\item{} \eqn{MPO_{heuristic} = 100 + 0.5 \times \max(0, HTL - 40)}{}.
\item{} \eqn{MPO_{safe} = LDL_{spl} - 5}{}.
\item{} \eqn{PTS_{safe\_limit} = 105 + 0.5 \times \max(0, HTL - 50)}{} (110 for Children).
\item{} \eqn{MPO_{final} = \min(120, MPO_{heuristic}, MPO_{safe}, PTS_{safe\_limit})}{}.

\end{itemize}


\end{enumerate}

\end{Section}
%
\begin{Author}
Gregory R. Warnes \email{greg@warnes.net}
\end{Author}
%
\begin{References}
ANSI S3.5-1997, "American National Standard Methods for 
Calculation of the Speech Intelligibility Index" American National
Standards Institute, New York. 

Other software programs for calculating SII are available from
\url{https://sii.to/html/programs.html}.

\end{References}
%
\begin{SeeAlso}
 SII Constants: \code{\LinkA{critical}{critical}}, and
\code{\LinkA{sic.critical}{sic.critical}} 
\end{SeeAlso}
%
\begin{Examples}
\begin{ExampleCode}

## Example C.1 from ANSI S3.5-1997 Annex C
sii.C1 <- sii(
              speech   = c(50.0, 40.0, 40.0, 30.0, 20.0,  0.0),
              noise    = c(70.0, 65.0, 45.0, 25.0,  1.0,-15.0),
              threshold= c( 0.0,  0.0,  0.0,  0.0,  0.0,  0.0),
              method="octave"
	      )
sii.C1                        # rounded to 2 digits by default
print(sii.C1$sii, digits=20)  # full precision
summary(sii.C1)               # full details
plot(sii.C1)                  # plot
plot(sii.C1, clinical=TRUE)   # clinical SPLogram plot
## The value given in the Standard is $0.504$.


	      
## Same calculation, but manually specify the frequencies
## and importance function, and use default for threshold

sii.C1 <- sii(
              speech   = c(50.0, 40.0, 40.0, 30.0, 20.0,  0.0),
              noise    = c(70.0, 65.0, 45.0, 25.0,  1.0,-15.0),
              method="octave",
              freq=c(250, 500, 1000, 2000, 4000, 8000),
	      importance=c(0.0617, 0.1671, 0.2373, 0.2648, 0.2142, 0.0549)
	      )
sii.C1	     

## Now perform the calculation using frequency weights for the Connected
## Speech Test (CST)
sii.CST <- sii(
               speech   = c(50.0, 40.0, 40.0, 30.0, 20.0,  0.0),
               noise    = c(70.0, 65.0, 45.0, 25.0,  1.0,-15.0),
               method="octave",
	       importance="CST"
	      )
round(sii.CST$table[,-c(5:7,13)],2)
sii.CST$sii

## Example C.2 from ANSI S3.5-1997 Annex C

sii.C2 <- sii(
              speech   = rep(54.0, 18),
              noise    = c(40.0, 30.0, 20.0, rep(0, 18-3) ),
              threshold= rep(0.0,  18),
              method="one-third"
              )
sii.C2$table[1:3,1:8]
sii.C2

## Interpolation example, for 8 frequencies using NU6 importance
## weight, default values for noise.
sii.left <- sii(
                speech="raised",
                threshold=c(25,25,30,35,45,45,55,60),
                freq=c(250, 500, 1000, 2000, 3000, 4000, 6000, 8000),
                method="critical",
                importance="NU6",
                interpolate=TRUE
                )
sii.left


\end{ExampleCode}
\end{Examples}
\printindex{}
\end{document}
